\documentclass[11pt]{article}

\usepackage[utf8]{inputenc}
\usepackage[T1]{fontenc}
\usepackage{lmodern}
\usepackage{geometry}
\usepackage{graphicx}
\usepackage{amsmath}
\usepackage{booktabs}
\usepackage{hyperref}

\geometry{margin=1in}

\title{A Minimal Generic Article Template for Paper Generation Testing}
\author{Template Library}
\date{\today}

\begin{document}

\maketitle

\begin{abstract}
This document serves as a minimal generic article template for testing paper generation workflows, structured text insertion, and LaTeX compilation in a template-independent setting.
\end{abstract}

\section{Introduction}
This generic article template is intended for baseline testing. It provides a simple environment for validating whether generated content can be inserted into a standard \LaTeX{} article structure without relying on conference- or journal-specific style files.

\section{Method}
We assume that generated text, equations, figures, tables, and citations may later be injected into this template for downstream testing.

\section{Example Equation}
A simple equation is shown below:
\begin{equation}
f(x) = x^2 + 2x + 1
\end{equation}

\section{Example Figure}
Figure~\ref{fig:placeholder} is a placeholder reference for future asset testing.

\begin{figure}[h]
\centering
\fbox{\rule{0pt}{1.5in}\rule{0.8\linewidth}{0pt}}
\caption{Placeholder figure for template testing.}
\label{fig:placeholder}
\end{figure}

\section{Example Table}
Table~\ref{tab:example} provides a minimal table example.

\begin{table}[h]
\centering
\begin{tabular}{ll}
\toprule
Item & Description \\
\midrule
Template & Generic article \\
Purpose & Baseline LaTeX testing \\
\bottomrule
\end{tabular}
\caption{A minimal example table.}
\label{tab:example}
\end{table}

\section{Conclusion}
This template provides a simple baseline for validating generated academic content and LaTeX compilation behavior.

\bibliographystyle{plain}
\bibliography{refs}

\end{document}
